\makeabstract
{
NCAM2 is required for gastrulation and body axis patterning during embryogenesis in Nematostella
}
{
Ilhan Ali(1), Anna Postnikova(1), Katerina Ragkousi(1)
}
{
\textsuperscript{1}Department of Biology, Amherst College, Amherst, MA, USA
}
{
Embryonic development is a fundamental biological process essential for life in all animals. Embryogenesis depends on the regulated expression of proteins that establish axial polarity and the adhesion of epithelial cells in the early embryo. The purpose of this study is to determine whether NCAM2 contributes to these processes, with the goal of uncovering a novel function for NCAM2 during embryonic development in the starlet sea anemone Nematostella vectensis.
}
{
Adult N. vectensis polyps were maintained at 17 ºC in dark conditions, and gametes were obtained by induced spawning. Embryos were fertilized and cultured at 21 °C. NCAM2, β-catenin, and GFP (control) were knocked down using shRNA electroporation. Embryos were fixed at late blastula (18 hpf) and gastrula (26 hpf) stages for phenotypic and molecular analyses. Protein localization was assessed by immunostaining for Cadherin-3 and phosphorylated myosin light chain, and staining with phalloidin and DAPI for F-actin and DNA. Gene expression of the transcription factors snailA and brachyury was assessed using HCR in situ hybridization. All samples were imaged using confocal microscopy (Zeiss LSM980). Structural predictions of NCAM2 were generated with AlphaFold3.
}
{
NCAM2 knockdown perturbed cell adhesion and gastrulation, resulting in embryos arresting at the blastula stage and cells protruding from the epithelium. It also caused the loss of basal localization of Cadherin-3. AlphaFold3 structural predictions suggested that NCAM2 forms a strong dimer between immunoglobulin and fibronectin domains of opposing proteins, consistent with a role in adhesion. HCR in situ hybridization of snailA and brachyury revealed that NCAM2 and β-catenin knockdown embryos exhibited similar phenotypes, including reduced brachyury expression at the blastopore, expanded snailA expression in the mesoderm, loss of body axis patterning, and failed gastrulation. Treatment of NCAM2 and β-catenin knockdown embryos with azakenpaullone, a selective inhibitor of GSK3β in the β-catenin destruction complex, rescued gastrulation and restored snailA and brachyury expression to near-baseline levels.
}
{
NCAM2 is required for epithelial integrity, cell adhesion, gastrulation, and axial body patterning in Nematostella embryos. Since Cadherin-3 remains highly expressed in NCAM2 knockdown embryos, the loss of adhesion alone does not explain the cell protrusions. NCAM2 knockdown also disrupts body axis patterning similar to β-catenin knockdown, with expansion of snailA and reduction of brachyury expression. Both phenotypes responded to azakenpaullone treatment. This suggests that NCAM2 may regulate signaling pathways that regulate cell migration and body axis patterning.
}

\newpage
